\partstart{The Thalean Carrier, Cover, and Relational Kernel}
\section{The Thalean Graph \texorpdfstring{$G_{60}\cong\mathrm{AT4val}[60,6]$}{G60 = AT4val[60,6]}}\label{sec:carrier}
\status{imported graph identification; independently reconstructible presentation}
Let $D$ be the dodecahedral graph and $\mathrm{CDC}(D)$ its canonical bipartite double cover. The carrier is represented by
% Original equation number(s) 2; expression unchanged.
\setcounter{equation}{1}
\begin{equation}
G_{60}\cong L(\operatorname{CDC}(D)),
\label{eq:carrierconstruction}
\end{equation}

The cover has vertices $(u,s)$ with $s\in\F_2$ and edges $(u,s)\sim(v,1-s)$ whenever $u\sim_Dv$. Its 40 vertices and 60 edges give a connected line graph with 60 vertices, 120 edges, degree four, and 40 triangles. The accompanying code reconstructs this presentation. The named census crosswalk is imported from the earlier manuscript \citep{cave2026becoming,psv2013} and is replayed against a published common-cover quotient in Section~\ref{sec:commoncover}.

The working quotient tower is
% Original equation number(s) 3; expression unchanged.
\setcounter{equation}{2}
\begin{equation}
G_{60}\xrightarrow{q_a}G_{30}\xrightarrow{q_b}G_{15},
\qquad G_{15}\cong L(\text{Petersen}).
\label{eq:tower}
\end{equation}

Cover-sheet exchange yields $G_{30}\cong L(D)$; the induced antipodal quotient yields $G_{15}\cong L(\mathrm{Petersen})$. Both neighborhood-bijective quotient maps are checked. Labels in this package are reproducible coordinates, not an identification with another project's station labels.

The two commuting deck involutions give
% Original equation number(s) 4; expression unchanged.
\setcounter{equation}{3}
\begin{equation}
V_4=\{1,a,b,ab\},\qquad a^2=b^2=1,\qquad ab=ba.
\end{equation}

Every regular fiber is a four-element $V_4$ torsor. Its three nonidentity translations give the three perfect matchings of an abstract $K_4$. The fiber itself has no carrier edges. The abstract triality $\Aut(V_4)\cong S_3$ permutes these pairings; its extension to a structured execution law is a separate property. Section~\ref{sec:mode} uses the registered orientation mechanism, not this abstract triality alone, to supply Mode.

Two incidence objects have different domains. The sector construction reports
% Original equation number(s) 5; expression unchanged.
\setcounter{equation}{4}
\begin{equation}
M_{\rm sec}M_{\rm sec}^{\mathsf T}=A_{15}^3+2A_{15}^2+2I.
\end{equation}

with $M_{\rm sec}$ of size $15\times30$. The later history matrix $B_{\rm hist}$ has size $30\times12$ and acts on twelve pentagonal histories. The matrices are kept distinct throughout \citep{cave2026becoming,historyconference}.

\begingroup
\renewcommand{\thesection}{2A}
\section{The AT4Val[60,4--6] Family}\label{sec:censusfamily}
\status{external literature and census identifications; numerical presentations and quotient maps replayed here}
The terminal index in $\mathrm{AT4Val}[60,j]$ is a census identifier, not a cycle length or an evolution parameter. The arc-transitive tetravalent census combines group-theoretic bounds and computer enumeration; its order-$60$ entries need not belong to a single construction family \citep{psv2013,psvbound,graphsym}. Here a specific three-member family is studied:
\begin{equation}
G_4:=\mathrm{AT4Val}[60,4],\qquad
G_5:=\mathrm{AT4Val}[60,5],\qquad
G_6:=\mathrm{AT4Val}[60,6]\cong G_{60}.
\tag{F1}\label{eq:familylabels}
\end{equation}
The subscripts $4,5,6$ in this subsection name census entries; they do not replace the order subscripts in $G_{15},G_{30},G_{60}$.

\subsection*{The independent literature anchor}
Minchenko and Wanless identify their graph $F_1$ as $I_{60,1}\cong\mathrm{AT4Val}[60][4]$ and give a Schreier coset construction with
\[
\Gamma=\mathbb Z_2\times\mathbb Z_2\times S_5,
\qquad H\cong D_8,\qquad [\Gamma:H]=480/8=60.
\]
Throughout this manuscript $D_8$ means the dihedral group of \emph{order eight}. The connection set and subgroup embedding, not the index alone, specify the graph. Despite the paper's title, this particular $F_1$ is a bipartite integral arc-transitive \emph{non-Cayley} graph. The relevant description is on printed page~396; the paper also obtains a bipartite $30$-vertex quotient $M_1\cong I_{30,3}$ in Section~4.6 \citep{minchenko2015}. That quotient is not identified here with either of the nonbipartite $30$-vertex graphs below. The publisher PDF is the primary mathematical source; the Monash publication record and the ResearchGate copy used in the initial investigation are separately listed in Appendix~\ref{app:familyreplay}.

\paragraph{Discovery provenance.}
According to the author's research record, the near-match $[60,4]$ surfaced incidentally during work on $[60,6]$. Scott Cave noticed the unfamiliar terminal index in the assistant's displayed research stream and requested follow-up. Inspection of the Minchenko--Wanless construction led to comparison with $[60,5]$ and then with the common cover. This records the route by which the family entered this project; it is not a claim that the catalogued graphs or their published constructions are new.

\subsection*{Census crosswalk and discriminating invariants}
The independent C4 construction records give
\[
G_4=\mathrm{C4}[60,11],\qquad
G_5=\mathrm{C4}[60,17],\qquad
G_6=\mathrm{C4}[60,13].
\]
All three have $60$ vertices, $120$ edges, valency four and reported full automorphism-group order $480$. Their local and global cycle geometries differ \citep{c46011,c46017,c46013}.
\begin{center}
\begin{tabular}{@{}lcccc@{}}
\toprule
Graph & Bipartite & Girth & Diameter & Consistent-cycle lengths\\
\midrule
$G_4$ & Yes & 6 & 6 & $6,6,10$\\
$G_5$ & No & 6 & 5 & $6,10,12$\\
$G_6$ & No & 3 & 6 & $3,10,12$\\
\bottomrule
\end{tabular}
\end{center}
Girth, diameter and bipartiteness are recomputed from the edge lists. The final column is imported from the census: it concerns consistent cycles, not every cycle of the graph. Matching lengths in this column do not by themselves identify cycle classes or give a transformation $G_4\to G_5\to G_6$.

\subsection*{Two different thirty-vertex intermediates}
Write $P$ for the Petersen graph and introduce
\begin{equation}
Q_{30}:=L(\mathrm{CDC}(P))\cong\mathrm{C4}[30,6].
\tag{F2}\label{eq:common30}
\end{equation}
The C4 name is also $\mathrm{Pr}_{10}(1,1,3,3)$; the construction record identifies it as $L(B(F_{10}))$, with $F_{10}=P$ and $B=\mathrm{CDC}$ \citep{c4306}. All three $G_i$ double-cover this $Q_{30}$, which in turn double-covers $L(P)$. In particular $G_4\cong\mathrm{CDC}(Q_{30})$. These relations are replayed below.

\begin{principle}[title={The two thirty-state graphs are not interchangeable}]
The original manuscript tower retains $G_{30}=L(D)$ in equation~\eqref{eq:tower}. The newly shared family base is $Q_{30}=L(\mathrm{CDC}(P))$. Exact graph isomorphism testing gives $Q_{30}\not\cong G_{30}$, although both have thirty vertices and cover $L(P)$. Thus the family supplies an additional route from $G_{60}$ to $G_{15}$, not a relabeling of the old route.
\end{principle}
The certified $G_{30}$-visible time register in Section~\ref{sec:time} remains attached to the original $G_{30}$. Nothing in the census comparison moves that visibility theorem to $Q_{30}$.
\endgroup
\setcounter{section}{2}

\begingroup
\renewcommand{\thesection}{2B}
\section{The Common Cover, Parity, and Rooted Frames}\label{sec:commoncover}
\status{published numerical input; exact finite reconstruction and elementary covering proofs}
Define the graph
\begin{equation}
\mathcal K_{120}:=\mathrm{C4}[120,38].
\tag{F3}\label{eq:commoncover}
\end{equation}
It has $120$ vertices and $240$ edges. Its published permutation generators produce a group of order $960$, equal to the full automorphism order reported by the census \citep{c412038,c412038forms}. The package verifies that every generator preserves the transcribed edge list and that the orbit of the published seed edge is exactly that list. Completeness of the automorphism group uses the external census order; the group closure and subsequent maps are computed here.

\begin{proposition}[Three central quotient faces]\label{prop:centralfaces}
In the published action on $\mathcal K_{120}$, the center is
\[
Z(\Aut(\mathcal K_{120}))=
\{1,\tau_4,\tau_5,\tau_6\}\cong V_4,
\qquad \tau_5\tau_6=\tau_4.
\]
Each nonidentity element acts freely on vertices. Its quotient is a locally bijective two-fold graph cover onto $G_i$:
\begin{equation}
\mathcal K_{120}/\langle\tau_i\rangle\cong G_i,
\qquad i\in\{4,5,6\}.
\tag{F4}\label{eq:threequotients}
\end{equation}
The element $\tau_4$ preserves the bipartition, while $\tau_5$ and $\tau_6$ exchange its parts. Quotienting by the full central $V_4$ gives $Q_{30}$.
\end{proposition}
\begin{proof}[Finite verification]
The script \path{scripts/verify_census_family.py} closes the four published generators, tests commutation with each generator, and obtains exactly four central permutations. It enumerates their vertex orbits and checks every local neighborhood map. The three quotient edge sets are compared by exact isomorphism with independently specified reference graphs: $\mathrm{CDC}(Q_{30})$, the published cyclic-voltage presentation of $\mathrm{C4}[60,17]$, and the inherited $L(\mathrm{CDC}(D))$ construction. The report contains the permutations, all quotient maps and the isomorphism witnesses. The product and color-action assertions are checked on all $120$ vertices.
\end{proof}

The construction records also identify
\begin{equation}
\mathrm{CDC}(G_5)\cong\mathcal K_{120}\cong\mathrm{CDC}(G_6).
\tag{F5}\label{eq:commondouble}
\end{equation}
This is consistent with the two color-exchanging quotient involutions \citep{c412038names}. The complete finite diagram is
\[
\begin{array}{ccccc}
&&\mathcal K_{120}&&\\
&\swarrow&\downarrow&\searrow&\\
G_4&&G_5&&G_6\cong G_{60}\\
&\searrow&\downarrow&\swarrow&\\
&&Q_{30}&&\\
&&\downarrow&&\\
&&G_{15}\cong L(P).&&
\end{array}
\]
The additional, original route $G_{60}\to G_{30}=L(D)\to G_{15}$ is retained separately.

The three $\tau_i$ are \emph{central}. Therefore no automorphism of this common graph conjugates $\tau_5$ to $\tau_6$. Different quotient choices are not thereby the two QR Modes of one event-presentation orbit. The correspondence with the abstract list $[1,a,b,ab]$ is a group-pattern comparison until a native cross-register map is given. In particular, the identity is a permutation, not itself a $120$-state history.

\subsection*{What the extra binary coordinate records}
For the Thalean branch, an explicit chart is
\begin{equation}
V(\mathcal K_{120})\cong V(G_{60})\times\F_2,
\qquad (v,\epsilon)\sim(w,\epsilon+1)
\ \Longleftrightarrow\ v\sim w.
\tag{F6}\label{eq:kparity}
\end{equation}
For an actual carrier walk $\gamma$ of length $n$, unique lifting from a specified initial sheet gives
\[
\epsilon_{\rm end}=\epsilon_{\rm start}+n\pmod2.
\]
Thus the canonical double retains \emph{carrier-walk parity}. It does not store the complete word or an integer winding. Once the initial lift is fixed there is exactly one lift of each base walk, not a fresh binary choice at every step.

The deck maps $a,b$ in equation~(4) are not carrier edges. Their canonical lifts $(v,\epsilon)\mapsto(a(v),\epsilon)$ and $(v,\epsilon)\mapsto(b(v),\epsilon)$ commute. Distinct ordered words $ab$ and $ba$ can still be retained as histories, but this canonical sheet coordinate does not distinguish their endpoint products. Nor does a common binary codomain identify carrier parity with an independently supplied $I/H$ field-transport signing; a common path domain and an explicit comparison of cocycles would be required.

Projection $(v,\epsilon)\mapsto v$ loses a distinction and is not by itself a complete information-preserving transduction. Retaining $(v,\epsilon)$ as propagated address plus receipt restores a bijective classical encoding. Describing the two lifts as \emph{immediate alternatives relative to a declared registration} is useful, but coherent superposition, a probability law and an outcome instrument require the additional structures of Section~\ref{sec:quantum}.

\subsection*{The eight-fold fiber and the two uses of \texorpdfstring{$D_8$}{D8}}
\begin{proposition}[Composite deck group and frame descent]\label{prop:deck8}
For the specified composite $\mathcal K_{120}\to G_{60}\to G_{15}$, the canonical lifts $\widetilde a,\widetilde b$ and the sheet flip $\kappa=\tau_6$ generate a free deck group
\begin{equation}
\langle\widetilde a,\widetilde b,\kappa\rangle\cong C_2^3.
\tag{F7}\label{eq:deck8}
\end{equation}
Its fifteen orbits have size eight and form the fibers of a locally bijective cover onto $L(P)$. The full published $960$-element group preserves these fibers. Its induced action has order $120$, kernel $C_2^3$, and a fiber stabilizer of order $64$. The stabilizer of a selected \emph{upstairs vertex} has order eight, is $D_8$, and maps faithfully onto the $D_8$ stabilizer of the corresponding $G_{15}$ vertex.
\end{proposition}
\begin{proof}
The three displayed generators are independent commuting involutions in the chart~\eqref{eq:kparity}. The inherited $V_4$ action is free on each four-fold fiber, and $\kappa$ changes only the sheet, so their product action is free on the eight-fold fiber. Composing the checked covers preserves local bijectivity. The script checks invariance of all fifteen blocks under the published generators and enumerates the induced $120$ permutations and its kernel. Root isotropy has element-order profile $1^1,2^5,4^2$; explicit generators satisfy $r^4=s^2=1$, $srs=r^{-1}$ and generate all eight elements. Their induced image on $G_{15}$ has order eight, hence the stabilizer map is injective and onto the corresponding rooted stabilizer.
\end{proof}

The positional identification is an identification of \emph{sets with action}:
\begin{equation}
V(G_{15})=E(P)\cong S_5/D_8^{\rm pos}.
\tag{F8}\label{eq:positionalcosets}
\end{equation}
In the two-subset presentation of $P$, an edge is a pair of disjoint two-subsets of five atoms. Fixing that unordered pair permits the two swaps within the subsets and their interchange, giving $D_8$. An ordered Petersen edge has a smaller stabilizer. Since $D_8^{\rm pos}$ is not normal in $S_5$, (F8) is a coset space, not a quotient group.

\begin{center}
\begin{tabularx}{\textwidth}{@{}p{.25\textwidth}X@{}}
\toprule Object & Action and status \\\midrule
$D_8^{\rm frame}$ & Root isotropy on $\mathcal K_{120}$, $G_{60}$ and $G_{15}$; the specified descent identifies these rooted actions. It is not the free eight-state deck group.\\
$D_8^{\rm hist}$ & A history/surface action in the imported $D_8\times C_2$ representation packet. Its domain is a history representation, not a rooted graph vertex \citep{historyframe2026}. A native WXYZTI or common-cover intertwiner is not supplied by that packet.\\
$C_2^3$ & The free eight-fold deck action of (F7), distinct from either use of $D_8$.\\
\bottomrule
\end{tabularx}
\end{center}
These are action types, not different abstract dihedral groups. An isomorphism of abstract groups does not identify their representations or physical roles. Likewise $960=|\Aut(\mathcal K_{120})|$ counts symmetries, not vertices; $\mathcal K_{120}$ still has $120$ vertices. Its central $V_4$ also rules out $\Aut(\mathcal K_{120})\cong S_5\times D_8$, whose center has order two.
\endgroup
\setcounter{section}{2}

\begingroup
\renewcommand{\thesection}{2C}
\section{A Twelve-Point Relational Kernel}\label{sec:relkernel}
\status{normalized construction imported from Project 41 Audits 028I--028J; relation algebra and symmetry replayed here}
A second finite object makes the local coupling question explicit without assuming a WXYZTI identification. Set
\begin{equation}
X_{12}=\mathbb Z_3\times V_4,
\qquad V_4\cong\F_2^2,
\qquad u=(1,0),\quad v=(0,1).
\tag{K1}\label{eq:kernelpoints}
\end{equation}
Its twelve points are $(i,x_1,x_2)$. The symbols $u,v$ label normalized coordinate directions, not new names for the historical deck permutations $a,b$. The source scripts obtain these coordinates from earlier finite constructions; the present replay begins with their explicit normalized laws \citep{rank8source,rank8group}.

\begin{definition}[Relational kernel]\label{def:relkernel}
The bounded \emph{Thalean relational kernel} $\mathscr K_{\rm rel}$ in this revision is the following eight-relation structure on $X_{12}$:
\begin{align}
T_t&=\{((i,x),(i,x+t)):i\in\Z_3,\ x\in\F_2^2\},
&&t\in\F_2^2,\notag\\
F_c&=\{((i,x),(i+1,y)):x_1+y_2=c\},
&&c\in\F_2,\tag{K2}\label{eq:kernelrelations}\\
R_c&=\{((i,x),(i-1,y)):x_2+y_1=c\},
&&c\in\F_2.\notag
\end{align}
Its adjacency algebra uses row-source, column-target matrices. This definition specifies a finite relational model; an execution grammar or a physical interpretation is additional data.
\end{definition}

\begin{proposition}[Rank-eight closure]\label{prop:rank8}
The relations $T_0,T_u,T_v,T_{u+v},F_0,F_1,R_0,R_1$ partition $X_{12}^2$. Each $T$ has valency one and each $F$ or $R$ has valency two. Transposition fixes the $T$ relations and exchanges $F_c$ with $R_c$. Every product of two adjacency matrices is an integer linear combination of the same eight matrices, with constant coefficients on each relation. The resulting homogeneous coherent configuration has rank eight and a noncommutative adjacency algebra.
\end{proposition}
\begin{proof}
Equal cell indices select exactly one translation; unequal indices select a direction and exactly one value of the indicated binary character. This proves the partition, valencies and transpose identities. The included exact replay evaluates all $64$ ordered matrix products on all $144$ ordered pairs and exports their constant intersection coefficients. Noncommutativity is already visible in
\begin{equation}
A_{T_u}A_{F_0}=A_{F_1},
\qquad A_{F_0}A_{T_u}=A_{F_0}.
\tag{K3}\label{eq:noncommutativewitness}
\end{equation}
Indeed, shifting the source $x_1$ reverses the forward character, while shifting the target $y_1$ leaves its $y_2$ constraint unchanged.
\end{proof}

This is a concrete informative dependence: under $F_0$, a source with $x_1=0$ admits precisely targets with $y_2=0$, whereas $x_1=1$ admits targets with $y_2=1$. An admitted target still has one unconstrained bit. Thus a relation can convey a distinction without being a deterministic permutation or already specifying a complete reversible event.

\begin{proposition}[Icosahedral fusion and color symmetry]\label{prop:fusion}
Fusing the eight matrices as
\begin{equation}
A_0=A_{T_0},\quad
A_1=A_{T_u}+A_{F_0}+A_{R_0},\quad
A_2=A_{T_v}+A_{F_1}+A_{R_1},\quad
A_3=A_{T_{u+v}}
\tag{K4}\label{eq:kernel-fusion}
\end{equation}
gives the distance-$0,1,2,3$ relations of the icosahedral graph, with intersection array $\{5,2,1;1,2,5\}$. The exact color-preserving automorphism group of the eight-relation structure is $C_2\times A_4$, of order $24$, transitive on twelve points with point stabilizer of order two. Its eight orbitals are precisely the eight relations.
\end{proposition}
\begin{proof}[Algebra and finite verification]
The replay identifies $A_1$ with the standard icosahedral graph and checks every distance relation and intersection count. A color-preserving permutation must preserve the three cells and their directed cyclic order. Preserving the named translations forces the form
\[
(i,x)\longmapsto(i+r,x+s_i),\qquad r\in\Z_3.
\]
Preservation of $F_0$ is exactly $(s_i)_1=(s_{i+1})_2$, so write $s_i=(b_{i+1},b_i)$ for an arbitrary $b\in\F_2^3$. There are $3\cdot 2^3=24$ such maps. Cell rotation cyclically permutes the three $b$ coordinates. The shift kernel splits into its diagonal line and even-parity plane; rotation fixes the former and cycles the three nonzero vectors of the latter. Hence the group is $C_2\times(V_4\rtimes C_3)\cong C_2\times A_4$. The orbit and orbital assertions are enumerated directly.
\end{proof}

The fusion forgets specified relational colors and retains an exact finite geometry. This is not a derivation of physical space. The full fused graph group has order $120$; the color group is self-normalizing with index five, yielding five conjugate refinements of this type. These assertions are also replayed, independently of a spacetime or polytope interpretation.

\subsection*{Points, relations and proposed physical readings}
The point count $12=3\cdot4$ and the cross-cell descriptors have different domains. A descriptor $(i,\text{forward/reverse},c)$ names one of \emph{twelve bins of directed pairs}; each bin contains eight pairs. These $96$ pairs, together with $48$ within-cell pairs, exhaust $144$ ordered pairs. They are not twelve new point states.

Consequently the proposed readings ``channel'', ``Mode'', and ``reciprocal role'' cannot be assigned to the three descriptor entries solely from $3\cdot2\cdot2=12$. In particular $F_c^{\mathsf T}=R_c$ is relation reversal, not yet the point action $\Mevt$ of Section~\ref{sec:mode}. Likewise the noncommutativity of relation matrices does not establish noncommuting primitive execution maps. The twelve kernel points, the twelve WXYZTI reciprocal pairs and the twelve Born-history pentagons remain distinct objects. Section~\ref{sec:wxyzti} specifies the crosswalk that would relate them.
\endgroup
\setcounter{section}{2}


\section{The Thalion}\label{sec:thalion}
\setcounter{theorem}{0}
\begin{definition}[Thalion]
A Thalion is a finite structured relational transducer whose local carrier is $G_{60}$, together with only those quotient, deck, transport, closure, history, and readout structures that have been explicitly supplied or derived.
\end{definition}

A convenient specification is
% Original equation number(s) 6; expression unchanged.
\setcounter{equation}{5}
\begin{equation}
\mathfrak T=(G_{60},q,\mathcal D,\{\mathsf T_e\},\mathcal R,\mathcal O),
\end{equation}

Here $q$ lists projections, $\mathcal D$ admissible contexts, $\mathsf T_e$ local operations, $\mathcal R$ the receipt/history interface, and $\mathcal O$ the readout family. A graph specifies possible adjacency; the admitted structure specifies the transducer.

The complete input includes prior local memory:
% Original equation number(s) 7; expression unchanged.
\setcounter{equation}{6}
\begin{equation}
\mathsf T_e:(x,m)\longmapsto(p,m'),
\label{eq:memorymap}
\end{equation}

The abbreviation $x\mapsto(p,r)$ is used only when $x$ already contains the complete input and $r$ the complete retained output. A Thalean transducer is a Thalion considered through these admitted input-output operations. Projected information may disappear from a readout while remaining recoverable in the complete output.
